\documentclass[UTF8,11pt]{ctexart}
\usepackage[a4paper,margin=24mm]{geometry}
\usepackage{amsmath,amssymb,amsthm,mathtools,mathrsfs}
\usepackage[all]{xy}
\usepackage{xurl,hyperref,enumitem}
\usepackage{tikz}
\usetikzlibrary{arrows.meta,cd,decorations.markings,calc,patterns}
\hypersetup{hidelinks,breaklinks=true}
\setlength{\emergencystretch}{3em}
\allowdisplaybreaks
\newtheorem*{theorem}{Theorem}
\newtheorem*{thm}{Theorem}
\newtheorem*{lemma}{Lemma}
\newtheorem*{proposition}{Proposition}
\newtheorem*{prop}{Proposition}
\newtheorem*{corollary}{Corollary}
\newtheorem*{cor}{Corollary}
\newtheorem*{claim}{Claim}
\theoremstyle{definition}
\newtheorem*{definition}{Definition}
\newtheorem*{defn}{Definition}
\newtheorem*{remark}{Remark}
\newtheorem*{example}{Example}
\newcommand{\R}{\mathbb R}
\newcommand{\C}{\mathbb C}
\newcommand{\Z}{\mathbb Z}
\newcommand{\Q}{\mathbb Q}
\newcommand{\N}{\mathbb N}
\newcommand{\F}{\mathbb F}
\newcommand{\D}{\mathbb D}
\newcommand{\bB}{\mathbb B}
\newcommand{\bC}{\mathbb C}
\newcommand{\bR}{\mathbb R}
\newcommand{\bZ}{\mathbb Z}
\newcommand{\bN}{\mathbb N}
\newcommand{\bQ}{\mathbb Q}
\newcommand{\bD}{\mathbb D}
\newcommand{\bF}{\mathbb F}
\newcommand{\CP}{\mathbb{CP}}
\newcommand{\RP}{\mathbb{RP}}
\newcommand{\sO}{\mathscr O}
\newcommand{\sI}{\mathscr I}
\newcommand{\sF}{\mathscr F}
\newcommand{\sG}{\mathscr G}
\newcommand{\sH}{\mathscr H}
\newcommand{\sR}{\mathscr R}
\newcommand{\sM}{\mathscr M}
\newcommand{\sV}{\mathscr V}
\newcommand{\sU}{\mathscr U}
\DeclareMathOperator{\Hom}{Hom}
\DeclareMathOperator{\Ext}{Ext}
\DeclareMathOperator{\Tor}{Tor}
\DeclareMathOperator{\Spec}{Spec}
\DeclareMathOperator{\Proj}{Proj}
\DeclareMathOperator{\supp}{supp}
\DeclareMathOperator{\im}{im}
\DeclareMathOperator{\id}{id}
\DeclareMathOperator{\Tr}{Tr}
\DeclareMathOperator{\tr}{tr}
\DeclareMathOperator{\rank}{rank}
\DeclareMathOperator{\ord}{ord}
\DeclareMathOperator{\dist}{dist}
\DeclareMathOperator{\End}{End}
\DeclareMathOperator{\Aut}{Aut}
\DeclareMathOperator{\Gal}{Gal}
\DeclareMathOperator{\vol}{vol}
\DeclareMathOperator{\Cl}{Cl}
\DeclareMathOperator{\coker}{coker}
\DeclareMathOperator{\codim}{codim}
\DeclareMathOperator{\divergence}{div}
\DeclareMathOperator{\Ric}{Ric}
\DeclareMathOperator{\grad}{grad}
\DeclareMathOperator{\Hess}{Hess}
\newcommand{\abs}[1]{\left\lvert#1\right\rvert}
\newcommand{\sabs}[1]{\lvert#1\rvert}
\newcommand{\babs}[1]{\bigl\lvert#1\bigr\rvert}
\newcommand{\Babs}[1]{\Bigl\lvert#1\Bigr\rvert}
\newcommand{\bbabs}[1]{\biggl\lvert#1\biggr\rvert}
\newcommand{\BBabs}[1]{\Biggl\lvert#1\Biggr\rvert}
\newcommand{\norm}[1]{\left\lVert#1\right\rVert}
\newcommand{\snorm}[1]{\lVert#1\rVert}
\newcommand{\bnorm}[1]{\bigl\lVert#1\bigr\rVert}
\newcommand{\Bnorm}[1]{\Bigl\lVert#1\Bigr\rVert}
\newcommand{\bbnorm}[1]{\biggl\lVert#1\biggr\rVert}
\newcommand{\BBnorm}[1]{\Biggl\lVert#1\Biggr\rVert}
\newcommand{\widebar}[1]{\overline{#1}}
\newcommand{\nicefrac}[2]{\frac{#1}{#2}}
\newcommand{\linnprod}[2]{\langle#1,#2\rangle}
\newcommand{\myindex}[1]{#1}
\newcommand{\glsadd}[1]{}
\newcommand{\avoidbreak}{}
\newcommand{\sourceref}[1]{\textnormal{[原文：\nolinkurl{#1}]}}
\newcommand{\thmref}[1]{\sourceref{#1}}
\newcommand{\lemmaref}[1]{\sourceref{#1}}
\newcommand{\propref}[1]{\sourceref{#1}}
\newcommand{\corref}[1]{\sourceref{#1}}
\newcommand{\defnref}[1]{\sourceref{#1}}
\newcommand{\exampleref}[1]{\sourceref{#1}}
\newcommand{\exerciseref}[1]{\sourceref{#1}}
\newcommand{\figureref}[1]{\sourceref{#1}}
\newcommand{\sectionref}[1]{\sourceref{#1}}
\newcommand{\appendixref}[1]{\sourceref{#1}}
\newcommand{\stacksref}[2]{\href{https://stacks.math.columbia.edu/tag/#1}{#2 [#1]}}


\newcommand{\E}{\mathbb E}
\newcommand{\Prob}{\mathbb P}
\newcommand{\Reff}{\mathcal R_{\mathrm{eff}}}
\newcommand{\eps}{\varepsilon}
\DeclareMathOperator{\diam}{diam}
\DeclareMathOperator{\Var}{Var}
\DeclareMathOperator{\Crit}{Crit}
\newcommand{\dd}{\,\mathrm d}


\begin{document}
\section*{074\quad 复仿射空间单射多项式映射的Ax--Grothendieck满射性}
\subsection*{来源与整理范围}
Terence Tao，\emph{Infinite fields, finite fields, and the Ax--Grothendieck theorem}，7 March 2009（网页注明3月8日修正）。Theorems 1--3及Section 1的Serre有限剩余域路线。
\par\url{https://terrytao.wordpress.com/2009/03/07/infinite-fields-finite-fields-and-the-ax-grothendieck-theorem/}
\par\textbf{恢复方式（整理者说明）：}从先前\nolinkurl{make074_075.py}执行记录恢复作者正文，重建独立排版；2026-09-28。
\subsection*{作者命题与人类证明}

\begin{theorem}[1, Ax--Grothendieck theorem in ${\bf C}$]
Let $P:{\bf C}^n\to{\bf C}^n$ be a polynomial. If $P$ is injective, then $P$ is bijective.
\end{theorem}
\paragraph{1. Proof via finite fields.}
The first observation is that the theorem is utterly trivial in the finite field case:
\begin{theorem}[2, Ax--Grothendieck theorem in $F$]
Let $F$ be a finite field, and let $P:F^n\to F^n$ be a polynomial. If $P$ is injective, then $P$ is bijective.
\end{theorem}
\begin{proof}
Any injection from a finite set to itself is necessarily bijective. (The hypothesis that $P$ is a polynomial is not needed at this stage, but becomes crucial later on.)
\end{proof}
Next, we pass from a finite field $F$ to its algebraic closure $\overline F$.
\begin{theorem}[3, Ax--Grothendieck theorem in $\overline F$]
Let $F$ be a finite field, let $\overline F$ be its algebraic closure, and let $P:\overline F^n\to\overline F^n$ be a polynomial. If $P$ is injective, then $P$ is bijective.
\end{theorem}
\begin{proof}
Our main tool here is Hilbert's nullstellensatz, which we interpret here as an assertion that if an algebraic problem is insoluble, then there exists a finitary algebraic obstruction that witnesses this lack of solution. Specifically, suppose for contradiction that we can find a polynomial $P:\overline F^n\to\overline F^n$ which is injective but not surjective. Injectivity of $P$ means that the algebraic system
\[P(x)=P(y);\qquad x\ne y\]
has no solution over the algebraically closed field $\overline F$; by the nullstellensatz, this implies that for each $j=1,\dots,n$, there must exist an algebraic identity of the form
\[(P(x)-P(y))\cdot Q_j(x,y)=(x_j-y_j)^r \tag{1}\]
for some $r\ge1$ and some polynomial $Q:\overline F^n\times\overline F^n\to\overline F^n$ that specifically witnesses this lack of solvability. Similarly, lack of surjectivity means the existence of an $z_0\in\overline F^n$ such that the algebraic system
\[P(x)=z_0\]
has no solution over the algebraically closed field $\overline F$; by another application of the nullstellensatz, there must exist an algebraic identity of the form
\[(P(x)-z_0)\cdot R(x)=1 \tag{2}\]
for some polynomial $R:\overline F^n\to\overline F^n$ that specifically witnesses this lack of solvability.

Fix $Q_j,z_0,R$ as above, and let $k$ be the subfield of $\overline F$ generated by $F$ and the coefficients of $P,Q_j,z_0,R$. Then we observe (thanks to our explicit witnesses (1), (2)) that the counterexample $P$ descends from $\overline F$ to $k$; $P$ is a polynomial from $k^n$ to $k^n$ which is injective but not surjective. But $k$ is finitely generated, and every element of $k$ is algebraic over the finite field $F$, thus $k$ is finite. But this contradicts Theorem 2.
\end{proof}
\begin{proof}[The complex case; the author's continuation]
The complex case ${\bf C}$ follows by a slight extension of the argument used to prove Theorem 3. Indeed, suppose for contradiction that there is a polynomial $P:{\bf C}^n\to{\bf C}^n$ which is injective but not surjective. As ${\bf C}$ is algebraically closed, we may invoke the nullstellensatz as before and find witnesses (1), (2) for some $Q,z_0,R$.

Now let $k={\bf Q}[\mathcal C]$ be the subfield of ${\bf C}$ generated by the rationals ${\bf Q}$ and the coefficients $\mathcal C$ of $P,Q,z_0,R$. Then we can descend the counterexample to $k$. This time, $k$ is not finite, but we can descend it to a finite field (and obtain the desired contradiction) by a number of methods.

One approach, which is the one taken by Serre, is to quotient the ring ${\bf Z}[\mathcal C]$ generated by the above coefficients by a maximal ideal, observing that this quotient is necessarily a finite field.
\end{proof}

\subsection*{整理说明与先修范围（非作者正文）}
按先前生成脚本恢复正文。按作者顺序收录有限域、有限域代数闭包以及复数域的Serre路线；作者随后给出的另一种超越基特化草图和Rudin路线不取入，也不分别计数。允许先修：Hilbert零点定理，以及有限生成$\Z$-代数的极大理想剩余域为有限域（Jacobson环的标准结果）。本题承担的核心工作是把单射性和非满射性同时编码为有限多项式恒等式，使两者都能经系数环商映射下降。$k=\Q[\mathcal C]$沿用作者作为生成子域的记号；证明分段标题与proof环境为编辑排版，正文未用AI补写。
\subsection*{可修改的简短标注（非作者正文）}
$c$：复多项式自映射的单射与满射关系。

$a$：零点定理；有限恒等式见证；有限生成系数环；极大理想剩余域；有限集计数。

\texttt{cross\_theory\_status = yes}。通过零点定理把无限域上的映射性质转换为可特化的代数恒等式，再返回有限集上的单射必满射矛盾。位置：式(1)、(2)和复数域收尾。

全部标注为非穷尽、可修改初稿。
\subsection*{来源许可与编辑归属}
作者About页Copyright etc.允许复制合理部分（如单篇文章）并注明来源。本题依该许可摘录：\url{https://terrytao.wordpress.com/about/}。
ChatGPT 加入题号、中文来源说明、整理范围和初稿标注；编辑说明与人类证明分列。
\end{document}
